\begin{table}[h!]
\label{tbl:riskassesmentexample}
\definecolor{headerblue}{RGB}{220, 225, 235}
\definecolor{riskyellow}{RGB}{255, 192, 0}

\begin{tabularx}{\textwidth}{|X|}
\hline
\rowcolor{headerblue}
\textbf{Notebooky N1} \\
\rowcolor{headerblue}
Dôvernosť: veľmi vysoká \\
\rowcolor{headerblue}
Integrita: vysoká \\
\rowcolor{headerblue}
Dostupnosť: veľmi vysoká \\
\hline
\end{tabularx}

\begin{tabularx}{\textwidth}{|X|X|}
\hline
Hrozba \newline
G 0.30 \textit{Neoprávnené používanie alebo správa zariadení a systémov}
&
Narušené základné hodnoty: \newline
dôvernosť, integrita, dostupnosť \\
\hline
\end{tabularx}

\begin{tabularx}{\textwidth}{|X|X|X|}
\hline
Frekvencia výskytu bez dodatočných bezpečnostných opatrení: často
&
Potenciálny dopad bez dodatočných bezpečnostných opatrení: značné
&
\rowcolor{riskyellow}
Riziko bez dodatočných bezpečnostných opatrení: vysoké \\
\hline
\end{tabularx}

\begin{tabularx}{\textwidth}{|X|}
\hline
\textbf{Popis:} \newline
Zamestnanci si pravidelne nosia notebooky na súťažné kolá, kde riadia registráciu a pomáhajú organizátorom podujatia a
z času na čas nakrátko odbehnú. Na notebookoch majú uložené rôzne údaje a materiály, častokrát aj správne súťažné
riešenia pre súťažne kolo, v prípade potreby tlače. Navyše, cez webový prehliadač pristupujú k intranetu organizácie,
kde sú vystavené mnohé neverejné dokumenty.
\newline\newline
\textbf{Vyhodnotenie:} \newline
Administrátori informačných systémov musia zaviesť povinné uzamykanie obrazovky, a to aj zavedením časového limitu pre automatické uzamknutie
zariadenia. Dodatočne by mali zaviesť šifrovanie diskov, aby predišli strate dát a úniku citlivých informácií v prípade narušenia integrity
zariadenia.
\hline
\end{tabularx}
\end{table}